%%%%%%%%%%%%%%%%%%%%%%%%%%%%%%%%%%%%%%%%%%%%%%%%%%%%%%%%%%%%%%%%
%%%% Standard Packages

\usepackage{amsthm,graphicx,xcolor,verbatim,bbm,amsmath,amsfonts,amssymb,graphicx,enumitem,bm,mathtools,booktabs}
\IfFileExists{dsfont.sty}{\usepackage{dsfont}}{}
\IfFileExists{newclude.sty}{\usepackage{newclude}}{}



%%%%%%%%%%%%%%%%%%%%%%%%%%%%%%%%%%%%%%%%%%%%%%%%%%%%%%%%%%%%%%%%
%%%% Additional Packages

% Tikz pictures
\usepackage{tikz}
\usetikzlibrary{arrows.meta}
\usepackage{amsmath}
\tikzset{
  bowl/.style = {line width=1.2pt, draw=gray!70},
  etaLess/.style = {-{Stealth[length=3.2pt]}, line width=1pt, draw=teal!70!black},
  etaMore/.style = {-{Stealth[length=3.2pt]}, line width=1pt, draw=orange!80!black},
}


% Ours imports
\IfFileExists{fullpage.sty}{\usepackage{fullpage}}{}

\IfFileExists{pgfplots.sty}{\usepackage{pgfplots}}{}
\usepackage{xcolor}
\definecolor{PythonGreen}{HTML}{2ca02c}
\usepackage{url}            % simple URL typesetting
\IfFileExists{nicefrac.sty}{\usepackage{nicefrac}}{}
\IfFileExists{caption.sty}{%
  \usepackage{caption}
  \captionsetup{font=small}
}{}
\IfFileExists{float.sty}{\usepackage{float}}{}
\IfFileExists{mdframed.sty}{\usepackage{mdframed}}{}
\IfFileExists{fontawesome5.sty}{\usepackage{fontawesome5}}{}
\IfFileExists{wrapfig.sty}{\usepackage{wrapfig}}{}
\usepackage[T1]{fontenc}    % To write the name of Jastrz\k{e}bski
\usepackage[utf8]{inputenc}

\IfFileExists{subcaption.sty}{\usepackage{subcaption}}{}
\IfFileExists{overpic.sty}{\usepackage{overpic}}{}



\usepackage{hyperref}
\PassOptionsToPackage{numbers,compress}{natbib}
\usepackage{natbib}



\hypersetup{colorlinks=true, linktocpage}

% \usepackage[capitalize,noabbrev]{cleveref}

\hypersetup{
    colorlinks,
    linkcolor={red!50!black},
    citecolor={blue!80!black},
    urlcolor={blue!80!black}
}




%%%%%%%%%%%%%%%%%%%%%%%%%%%%%%%%%%%%%%%%%%%%%%%%%%%%%%%%%%%%%%%%
%%%% Macros

\newcommand{\pier}[1]{{\color{red}Pier: #1}}
\newcommand{\red}[1]{{\color{red} #1}}
\newcommand{\blue}[1]{{\color{blue} #1}}
\newcommand{\qandq}{\quad \text{and} \quad}
\newcommand{\lr}[1]{ \left( #1 \right) }
\newcommand{\norm}[1]{ \left\| #1 \right\| }
\newcommand{\mednorm}[1]{\| #1\| }
% \newcommand{\abs}[1]{ \left| #1 \right| }
% \newcommand{\ceil}[1]{ \left\lceil #1 \right\rceil }
% \newcommand{\floor}[1]{ \left\lfloor #1 \right\rfloor }
% \newcommand{\indicator}{\mathds{1}}


% \DeclareMathOperator*{\argmax}{arg\,max}
% \DeclareMathOperator*{\argmin}{arg\,min}
% \DeclareMathOperator{\Relu}{ReLU}
% \DeclareMathOperator{\reg}{reg}
% \DeclareMathOperator{\loss}{loss}
\DeclareMathOperator{\var}{\mathbb{V}ar}
\DeclareMathOperator{\Cov}{Cov}
% \DeclareMathOperator{\tr}{tr}
% \DeclareMathOperator{\dist}{dist}
% \DeclareMathOperator{\sgn}{sign}

\newcommand{\Borel}{\mathcal{B}}
\newcommand{\Borelmeasure}{\operatorname{Borel}}
\newcommand{\E}{\mathbb{E}}
\renewcommand{\P}{\mathbb{P}}
% \newcommand{\Q}{\mathbb{Q}}
\newcommand{\R}{\mathbb{R}}
% \newcommand{\S}{\mathbb{S}}
\newcommand{\N}{\mathbb{N}}
\newcommand{\Sym}{\mathbb{S}}
% \newcommand{\Z}{\mathbb{Z}}
\newcommand{\cV}{\mathcal{V}}
\newcommand{\smallsum}{\textstyle\sum}
% \newcommand{\textint}{\textstyle\int}
\newcommand{\scp}[3]{\langle#1,#2\rangle_{#3}}
\newcommand{\overbar}[1]{\mkern 1.5mu\overline{\mkern-1.5mu#1\mkern-1.5mu}\mkern 1.5mu}

\newcommand{\limitarrow}[1]{%
\parbox{#1}{\tikz{\draw[->](0,0)--(#1,0);}}}





\newcommand{\A}{\mathbf{A}}
\renewcommand{\a}[0]{\mathbf{a}}
\renewcommand{\b}[0]{\mathbf{b}}
\renewcommand{\c}[0]{\mathbf{c}}
\newcommand{\atb}[0]{\mathbf{a}^\top \mathbf{b}}
\newcommand{\x}[0]{\mathbf{x}}
\newcommand{\y}[0]{\mathbf{y}}

\DeclareMathOperator{\relu}{\mathrm{ReLU}}
\DeclareMathOperator{\tr}{\mathrm{Tr}}
\DeclareMathOperator*{\argmin}{arg\,min}
\DeclareMathOperator*{\argmax}{arg\,max} 

\newcommand{\scale}[0]{\boldsymbol{\lambda}}
\newcommand{\residual}[0]{\boldsymbol{\varepsilon}}



%%%% Blake Macros

\DeclarePairedDelimiter{\abs}{\lvert}{\rvert} %
\DeclarePairedDelimiter{\brk}{[}{]}
\DeclarePairedDelimiter{\crl}{\{}{\}}
\DeclarePairedDelimiter{\prn}{(}{)}
\DeclarePairedDelimiter{\nrm}{\|}{\|}
\DeclarePairedDelimiter{\ceil}{\lceil}{\rceil}
\DeclarePairedDelimiter{\floor}{\lfloor}{\rfloor}
\newcommand{\inner}[2]{\left\langle #1,\, #2 \right\rangle}






%%%%%%%%%%%%%%%%%%%%%%%%%%%%%%%%%%%%%%%%%%%%%%%%%%%%%%%%%%%%%\
%%%% Define commands

\theoremstyle{plain}
\newtheorem{theorem}{Theorem}%[section]
% \newtheorem*{dreamtheorem}{Conjecture}%[dreamtheorem]
% \newtheorem*{theorem*}{Theorem}%[section]
% \newtheorem*{theorem_1}{Theorem \ref{theo:1}}
% \newtheorem*{theorem_0}{Theorem \ref{theo:0}}

\newtheorem{lemma}{Lemma}
\newtheorem{proposition}{Proposition}
% \newtheorem{cor}[theorem]{Corollary}
% \newtheorem{example}[theorem]{Example}
% \newtheorem{setting}{Setting}
% \newtheorem{conjecture}[theorem]{Conjecture}
% \newtheorem{dconjecture}[theorem]{Draftconjecture}
\theoremstyle{remark}
\newtheorem{remark}{Remark}
% \newtheorem{remark}[theorem]{Remark}
% \newtheorem{condition}[theorem]{Condition}
% % \newtheorem*{condition*}[theorem]{Condition}
\theoremstyle{definition}
\newtheorem{definition}{Definition}
\newtheorem{assumption}{Assumption}
\newtheorem{hypothesis}{Hypothesis}
\newtheorem{question}{Question}
\newtheorem{example}{Example}

% Starred versions for restating without incrementing numbers
\newtheorem*{theorem*}{Theorem}
\newtheorem*{lemma*}{Lemma}
\newtheorem*{proposition*}{Proposition}
\newtheorem*{definition*}{Definition}

% Define any custom commands that you want to use.
% For example, highlight notes for future edits to the thesis
%\newcommand{\todo}[1]{\textbf{\emph{TODO:}#1}}


%%%%%%%%%%%%%%%%%%%%%%%%%%%%%%%%
% MACROS (lightweight, optional)
%%%%%%%%%%%%%%%%%%%%%%%%%%%%%%%%
\newcommand{\EoS}{\textsc{EoS}}
\newcommand{\EoSS}{\textsc{EoSS}}
\newcommand{\CE}{\textsc{CE}}
\newcommand{\MSC}{\textsc{MSC}}
\newcommand{\MSE}{\textsc{MSE}}
\newcommand{\GD}{\textsc{GD}}
\newcommand{\SGD}{\textsc{SGD}}
\newcommand{\AdamW}{\textsc{AdamW}}



\IfFileExists{tcolorbox.sty}{%
  \usepackage[most]{tcolorbox}
  \newtcolorbox{summarybox}{
    % colback=blue!5,          % light background
    colframe=black!70,         % border color
    boxrule=2pt,               % border thickness
    arc=4pt,                   % rounded corners
    left=2pt,
    right=2pt,
    top=3pt,
    bottom=3pt,
    halign=center,
  }
}{%
  \newenvironment{summarybox}{\begin{center}\begin{minipage}{0.95\linewidth}}{\end{minipage}\end{center}}
}



%%%% Possibly to remove
\setlength{\parindent}{0pt}
\setlength{\parskip}{5.5pt}
